\ifdefined\gkver\else\def\gkver{student}\fi
\documentclass[\gkver]{gaokaozhenti}
\begin{document}

\chapter{2013年浙江理综}
%% paperType: 理综物理部分

\begin{enumerate}
\item
%% number: 14
%% paperName: 2013年浙江理综第 14 题 4 分
%% typeId: 1
%% type: 单选题
%% chapter: 光学
%% point: 光的电磁说
%% method: 无
%% score: 4
%% degree: 850
%% duplicateId: 0
%% body:
关于生活中遇到的各种波，下列说法正确的是\xzanswer{B}

\fourchoices[answer=B]
{电磁波可以传递信息，声波不能传递信息}
{手机在通话时涉及的波既有电磁波又有声波}
{太阳光中的可见光和医院“B 超”中的超声波传递速度相同}
{遥控器发出的红外线波长和医院 CT 中的 X 射线波长相同}

%% answer: B
%% memo:
\memoanswer{
本题原卷及材料中未提供详解，待补充。
}

\item
%% number: 15
%% paperName: 2013年浙江理综第 15 题 4 分
%% typeId: 1
%% type: 单选题
%% chapter: 电磁感应
%% point: 电磁感应中图象
%% method: 无
%% score: 4
%% degree: 850
%% duplicateId: 0
%% body:
磁卡的磁条中有用于存储信息的磁极方向不同的磁化区，刷卡器中有检测线圈，当以速度 $v_{0}$ 刷卡时，在线圈中产生感应电动势。其 $E-t$ 关系如右图所示。如果只将刷卡速度改为 $v_{0}/2$，线圈中的 $E-t$ 关系可能是\xzanswer{D}

\onepicture[w=8cm, align=right]{figs/15.png}

\fourchoices[answer=D, ispicture=true, h=1.6cm]
{figs/15a.png}
{figs/15b.png}
{figs/15c.png}
{figs/15d.png}

%% answer: D
%% memo:
\memoanswer{
本题原卷及材料中未提供详解，待补充。
}

\item
%% number: 16
%% paperName: 2013年浙江理综第 16 题 4 分
%% typeId: 1
%% type: 单选题
%% chapter: 光学
%% point: 光的折射
%% method: 无
%% score: 4
%% degree: 850
%% duplicateId: 0
%% body:
与通常观察到的月全食不同，小虎同学在2012年12月10日晚观看月全食时，看到整个月亮是暗红的。小虎画了月全食的示意图，并提出了如下猜想，其中最为合理的是\xzanswer{C}

\onepicture[w=4.5cm, align=right]{figs/16.png}

\fourchoices[answer=C]
{地球上有人用红色激光照射月球}
{太阳照射到地球的红光反射到月球}
{太阳光中的红光经地球大气层折射到月球}
{太阳光中的红光在月球表面形成干涉条纹}

%% answer: C
%% memo:
\memoanswer{
本题原卷及材料中未提供详解，待补充。
}

\item
%% number: 17
%% paperName: 2013年浙江理综第 17 题 4 分
%% typeId: 1
%% type: 单选题
%% chapter: 运动和力
%% point: 牛顿定律的应用
%% method: 无
%% score: 4
%% degree: 650
%% duplicateId: 0
%% body:
如图所示，水平板上有质量 $m=1.0\Ukg$ 的物块，受到随时间 $t$ 变化的水平拉力 $F$ 作用，用力传感器测出相应时刻物块所受摩擦力 $F_{f}$ 的大小。取重力加速度 $g=10\Umsq$。下列判断正确的是\xzanswer{D}

\onepicture[w=9cm, align=right]{figs/17.png}

\fourchoices[answer=D]
{$5\Us$ 内拉力对物块做功为零}
{$4\Us$ 末物块所受合力大小为 $4.0\UN$}
{物块与木板之间的动摩擦因数为 0.4}
{$6\Us\sim9\Us$ 内物块的加速度的大小为 $2.0\Umsq$}

%% answer: D
%% memo:
\memoanswer{
本题原卷及材料中未提供详解，待补充。
}

\item
%% number: 18
%% paperName: 2013年浙江理综第 18 题 6 分
%% typeId: 2
%% type: 多选题
%% chapter: 曲线运动
%% point: 人造地球卫星
%% method: 无
%% score: 6
%% degree: 650
%% duplicateId: 0
%% body:
如图所示，三颗质量均为 $m$ 的地球同步卫星等间隔分布在半径为 $r$ 的圆轨道上，设地球质量为 $M$，半径为 $R$。下列说法正确的是\xzanswer{BC}

\onepicture[w=4.5cm, align=right]{figs/18.png}

\fourchoices[answer=BC]
{地球对一颗卫星的引力大小为 $\frac{GMm}{(r-R)^{2}}$}
{一颗卫星对地球的引力大小为 $\frac{GMm}{r^{2}}$}
{两颗卫星之间的引力大小为 $\frac{Gm^{2}}{3r^{2}}$}
{三颗卫星对地球引力的合力大小为 $\frac{3GMm}{r^{2}}$}

%% answer: BC
%% memo:
\memoanswer{
本题原卷及材料中未提供详解，待补充。
}

\item
%% number: 19
%% paperName: 2013年浙江理综第 19 题 6 分
%% typeId: 2
%% type: 多选题
%% chapter: 运动和力
%% point: 牛顿定律的应用
%% method: 无
%% score: 6
%% degree: 850
%% duplicateId: 0
%% body:
如图所示，总质量为 $460\Ukg$ 的热气球，从地面刚开始竖直上升时的加速度为 $0.5\Umsq$，当热气球上升到 $180\Um$ 时，以 $5\Ums$ 的速度向上匀速运动。若离开地面后热气球所受浮力保持不变，上升过程中热气球总质量不变，重力加速度 $g=10\Umsq$。关于热气球，下列说法正确的是\xzanswer{AD}

\onepicture[w=3cm, align=right]{figs/19.png}

\fourchoices[answer=AD]
{所受浮力大小为 $4830\UN$}
{加速上升过程中所受空气阻力保持不变}
{从地面开始上升 $10\Us$ 后的速度大小为 $5\Ums$}
{以 $5\Ums$ 匀速上升时所受空气阻力大小为 $230\UN$}

%% answer: AD
%% memo:
\memoanswer{
本题原卷及材料中未提供详解，待补充。
}

\item
%% number: 20
%% paperName: 2013年浙江理综第 20 题 6 分
%% typeId: 2
%% type: 多选题
%% chapter: 磁场
%% point: 带电粒子在磁场中的运动
%% method: 无
%% score: 6
%% degree: 650
%% duplicateId: 0
%% body:
在半导体离子注入工艺中，初速度可忽略的离子 $\ce{P+}$ 和 $\ce{P^{3+}}$，经电压为 $U$ 的电场加速后，垂直进入磁感应强度大小为 $B$、方向垂直纸面向里，有一定的宽度的匀强磁场区域，如图所示。已知离子 $\ce{P+}$ 在磁场中转过 $\theta=30^{\circ}$ 后从磁场右边界射出。在电场和磁场中运动时，离子 $\ce{P+}$ 和 $\ce{P^{3+}}$\xzanswer{BCD}

\onepicture[w=4cm, align=right]{figs/20.png}

\fourchoices[answer=BCD]
{在电场中的加速度之比为 $1:1$}
{在磁场中运动的半径之比为 $3:1$}
{在磁场中转过的角度之比为 $1:2$}
{离开电场区域时的动能之比为 $1:3$}

%% answer: BCD
%% memo:
\memoanswer{
本题原卷及材料中未提供详解，待补充。
}

\item
%% number: 21
%% paperName: 2013年浙江理综第 21 题 10 分
%% typeId: 4
%% type: 实验题
%% chapter: 直线运动
%% point: 打点计时器公式
%% method: 无
%% score: 10
%% degree: 650
%% duplicateId: 0
%% body:
如图所示，装置甲中挂有小桶的细线绕过定滑轮，固定在小车上；装置乙中橡皮筋的一端固定在导轨的左端，另一端系在小车上。一同学用装置甲和乙分别进行实验，经正确操作获得两条纸带①和②，纸带上的 a、b、c $\ldots\ldots$ 均为打点计时器打出的点。
\begin{enumerate}
	\item
	任选一条纸带读出 b、c 两点间的距离为 \_\_\_\_\_\_\_；
	\item
	任选一条纸带求出 c、e 两点间的平均速度大小为 \_\_\_\_\_\_\_，纸带①和②上 c、e 两点间的平均速度 $\overline{v}$①\_\_\_\_\_\_ $\overline{v}$②（填“大于”、“等于”或“小于”）；
	\item
	图中 \_\_\_\_\_\_\_（填选项）

	\fourchoices[answer=C]
	{两条纸带均为用装置甲实验所得}
	{两条纸带均为用装置乙实验所得}
	{纸带①为用装置甲实验所得，纸带②为用装置乙实验所得}
	{纸带①为用装置乙实验所得，纸带②为用装置甲实验所得}
\end{enumerate}

\onepicture[w=11cm, align=right]{figs/21.png}

%% answer:
\jdanswer{
\begin{enumerate}
	\item
	$2.10\Ucm$ 或 $2.40\Ucm$（$\pm0.05\Ucm$，有效数字位数正确）

	\item
	$1.13\Ums$ 或 $1.25\Ums$（$\pm0.05\Ums$，有效数字位数不作要求），小于

	\item
	C

\end{enumerate}
}
%% memo:
\memoanswer{
本题原卷及材料中未提供详解，待补充。
}

\item
%% number: 22
%% paperName: 2013年浙江理综第 22 题 10 分
%% typeId: 4
%% type: 实验题
%% chapter: 电路
%% point: 测电动势与内阻
%% method: 无
%% score: 10
%% degree: 650
%% duplicateId: 0
%% body:
采用如图所示的电路“测定电池的电动势和内阻”。
\begin{enumerate}
	\item
	除了选用照片中的部分器材外，\_\_\_\_\_\_\_\_\_\_\_\_\_（填选项）

	\fourchoices[answer=A]
	{还需要电压表}
	{还需要电流表}
	{还需要学生电源}
	{不再需要任何器材}

	\item
	测量所得数据如下：

	\begin{center}
	\begin{tabular}{|c|c|c|c|c|c|c|}
	\hline
	\diagbox{物理量}{测量次数} & 1 & 2 & 3 & 4 & 5 & 6 \\ \hline
	$R/\UO$ & 1.2 & 1.0 & 0.8 & 0.6 & 0.4 & 0.2 \\ \hline
	$I/\UA$ & 0.60 & 0.70 & 0.80 & 0.89 & 1.00 & 1.20 \\ \hline
	$U/\UV$ & 0.90 & 0.78 & 0.74 & 0.67 & 0.62 & 0.43 \\ \hline
	\end{tabular}
	\end{center}

	用作图法求得电池的内阻 $r=$\_\_\_\_\_\_\_\_；

	\item
	根据第5组所测得的实验数据，求得电流表内阻 $R_{A}=$\_\_\_\_\_\_。
\end{enumerate}

\onepicture[w=10cm, align=right]{figs/22.png}

%% answer:
\jdanswer{
\begin{enumerate}
	\item
	A

	\item
	见下图，$r=(0.75\pm0.10)\UO$

	\item
	$0.22\UO$

\end{enumerate}

\onepicture[w=8cm]{figs/22_an.jpg}
}
%% memo:
\memoanswer{
本题原卷及材料中未提供详解，待补充。
}

\item
%% number: 23
%% paperName: 2013年浙江理综第 23 题 18 分
%% typeId: 6
%% type: 计算题
%% chapter: 曲线运动
%% point: 平抛运动
%% method: 无
%% score: 18
%% degree: 650
%% duplicateId: 0
%% body:
山谷中有三块石头和一根不可伸长的轻质青藤，其示意图如下。图中A、B、C、D均为石头的边缘点，O为青藤的固定点，$h_{1}=1.8\Um$，$h_{2}=4.0\Um$，$x_{1}=4.8\Um$，$x_{2}=8.0\Um$。开始时，质量分别为 $M=10\Ukg$ 和 $m=2\Ukg$ 的大、小两只滇金丝猴分别位于左边和中间的石头上，当大猴发现小猴将受到伤害时，迅速从左边石头 A 点起水平跳到中间石头，大猴抱起小猴跑到 C 点，抓住青藤下端荡到右边石头上的 D 点，此时速度恰好为零。运动过程中猴子均看成质点，空气阻力不计，重力加速度 $g=10\Umsq$。求：
\begin{enumerate}
	\item
	大猴从 A 点水平跳离时速度的最小值；
	\item
	猴子抓住青藤荡起时的速度大小；
	\item
	猴子荡起时，青藤对猴子的拉力大小。
\end{enumerate}

\onepicture[w=7cm, align=right]{figs/23.png}

%% answer:
\jdanswer{
\begin{enumerate}
	\item
	$3\Ums$

	\item
	$\sqrt{80}\Ums\approx9\Ums$

	\item
	$216\UN$

\end{enumerate}
}
%% memo:
\memoanswer{
（1）设猴子从 A 点水平跳离时速度的最小值为 $v_{\min}$，根据平抛运动规律，有

$h_{1}=\frac{1}{2}gt^{2}$①，

$x_{1}=v_{\min}t$②，

联立①、②式得

$v_{\min}=3\Ums$③。

（2）猴子抓住青藤后的运动过程中机械能守恒，设荡起时速度为 $v_{c}$，有

$(M+m)gh_{2}=\frac{1}{2}(M+m)v_{c}^{2}$④，

$v_{c}=\sqrt{2gh_{2}}=\sqrt{80}\Ums\approx9\Ums$⑤。

（3）设拉力为 $F_{T}$，青藤的长度为 $L$，对最低点，由牛顿第二定律得

$F_{T}-(M+m)g=(M+m)\frac{v_{c}^{2}}{L}$⑥，

由几何关系

$(L-h_{2})^{2}+x_{2}^{2}=L^{2}$⑦，

得：$L=10\Um$⑧。

综合⑤、⑥、⑧式并代入数据解得：

$F_{T}=(M+m)g+(M+m)\frac{v_{c}^{2}}{L}=216\UN$。
}

\item
%% number: 24
%% paperName: 2013年浙江理综第 24 题 20 分
%% typeId: 6
%% type: 计算题
%% chapter: 电场
%% point: 带电粒子的圆周运动
%% method: 无
%% score: 20
%% degree: 650
%% duplicateId: 0
%% body:
“电子能量分析器”主要由处于真空中的电子偏转器和探测板组成。偏转器是由两个相互绝缘、半径分别为 $R_{A}$ 和 $R_{B}$ 的同心金属半球面 A 和 B 构成，A、B 为电势值不等的等势面，其过球心的截面如图所示。一束电荷量为 $e$、质量为 $m$ 的电子以不同的动能从偏转器左端 M 的正中间小孔垂直入射，进入偏转电场区域，最后到达偏转器右端的探测板 N，其中动能为 $E_{k0}$ 的电子沿等势面 C 做匀速圆周运动到达 N 板的正中间。忽略电场的边缘效应。
\begin{enumerate}
	\item
	判断球面 A、B 的电势高低，并说明理由；
	\item
	求等势面 C 所在处电场强度 $E$ 的大小；
	\item
	若半球面 A、B 和等势面 C 的电势分别为 $\varphi_{A}$、$\varphi_{B}$ 和 $\varphi_{C}$，则到达 N 板左、右边缘处的电子，经过偏转电场前、后的动能改变量 $\Delta E_{k\text{左}}$ 和 $\Delta E_{k\text{右}}$ 分别为多少？
	\item
	比较 $|\Delta E_{k\text{左}}|$ 和 $|\Delta E_{k\text{右}}|$ 的大小，并说明理由。
\end{enumerate}

\onepicture[w=6cm, align=right]{figs/24.png}

%% answer:
\jdanswer{
\begin{enumerate}
	\item
	B 板电势高于 A 板；

	\item
	$\frac{4E_{k0}}{e(R_{A}+R_{B})}$

	\item
	$\Delta E_{k\text{左}}=e(\varphi_{B}-\varphi_{C})$，$\Delta E_{k\text{右}}=e(\varphi_{A}-\varphi_{C})$

	\item
	$|\Delta E_{k\text{左}}|>|\Delta E_{k\text{右}}|$

\end{enumerate}
}
%% memo:
\memoanswer{
（1）电子（带负电）做圆周运动，电场力方向指向球心，电场方向从 B 指向 A，B 板电势高于 A 板。

（2）据题意，电子在电场力作用下做圆周运动，考虑轨道上的电场强度 $E$ 大小相同，有：

$eE=m\frac{v^{2}}{R}$，

$E_{k0}=\frac{1}{2}mv^{2}$，

$R=\frac{R_{A}+R_{B}}{2}$，

联立解得：

$E=\frac{2E_{k0}}{eR}=\frac{4E_{k0}}{e(R_{A}+R_{B})}$。

（3）电子运动时只有电场力做功，根据动能定理，有

$\Delta E_{k}=qU$，

对到达 N 板左边缘的电子，电场力做正功，动能增加，有

$\Delta E_{k\text{左}}=e(\varphi_{B}-\varphi_{C})$。

对到达 N 板右边缘的电子，电场力做负功，动能减少，有

$\Delta E_{k\text{右}}=e(\varphi_{A}-\varphi_{C})$。

（4）根据电场线特点，等势面 B 与 C 之间的电场强度大于 C 与 A 之间的电场强度，考虑到等势面间距相等，有

$|\varphi_{B}-\varphi_{C}|>|\varphi_{A}-\varphi_{C}|$，

即

$|\Delta E_{k\text{左}}|>|\Delta E_{k\text{右}}|$。
}

\item
%% number: 25
%% paperName: 2013年浙江理综第 25 题 22 分
%% typeId: 6
%% type: 计算题
%% chapter: 磁场
%% point: 安培力
%% method: 无
%% score: 22
%% degree: 450
%% duplicateId: 0
%% body:
为了降低潜艇噪音，提高其前进速度，可用电磁推进器替代螺旋桨。潜艇下方有左、右两组推进器，每组由6个相同的、用绝缘材料制成的直线通道推进器构成，其原理示意图如下。在直线通道内充满电阻率 $\rho=0.2\UOm$ 的海水，通道中 $a\times b\times c=0.3\Um\times0.4\Um\times0.3\Um$ 的空间内，存在由超导线圈产生的匀强磁场，其磁感应强度 $B=6.4\UT$、方向垂直通道侧面向外。磁场区域上、下方各有 $a\times b=0.3\Um\times0.4\Um$ 的金属板M、N，当其与推进器专用直流电源相连后，在两板之间的海水中产生了从N到M，大小恒为 $I=1.0\times10^{3}\UA$ 的电流，设电流只存在于磁场区域。不计电源内阻及导线电阻，海水密度 $\rho_{m}=1.0\times10^{3}\Ukgmc$。
\begin{enumerate}
	\item
	求一个直线通道推进器内磁场对通电海水的作用力大小，并判断其方向 ；
	\item
	在不改变潜艇结构的前提下，简述潜艇如何转弯？如何“倒车”？
	\item
	当潜艇以恒定速度 $v_{0}=30\Ums$ 前进时，海水在出口处相对于推进器的速度 $v=34\Ums$，思考专用直流电源所提供的电功率如何分配，求出相应功率的大小。
\end{enumerate}

\onepicture[w=11cm, align=right]{figs/25.png}

%% answer:
\jdanswer{
\begin{enumerate}
	\item
	$1.92\times10^{3}\UN$，方向向右

	\item
	考虑到潜艇下方有左、右两组推进器，可以开启或关闭不同个数的左、右两侧的直线通道推进器，实施转弯。 改变电流方向，或改变磁场方向，可以改变海水所受磁场力的方向，根据牛顿第三定律，使潜艇“倒车”。

	\item
	$4.6\times10^{4}\UW$

\end{enumerate}
}
%% memo:
\memoanswer{
（1）将通电海水看成导体，所受磁场力

$F=IBL$，

代入数据得：$F=IBc=1.0\times10^{3}\times6.4\times0.3\UN=1.92\times10^{3}\UN$。

用左手定则判断磁场对海水作用力方向向右（或与海水出口方向相同）。

（2）考虑到潜艇下方有左、右两组推进器，可以开启或关闭不同个数的左、右两侧的直线通道推进器，实施转弯。

改变电流方向，或改变磁场方向，可以改变海水所受磁场力的方向，根据牛顿第三定律，使潜艇“倒车”。

（3）电源提供的电功率中的第一部分：牵引功率

$P_{1}=F_{\text{牵}}v_{0}$，

根据牛顿第三定律：

$F_{\text{牵}}=12IBL$，

当 $v_{0}=30\Ums$ 时，代入数据得：

$P_{1}=F_{\text{牵}}v_{0}=12\times1.92\times10^{3}\times30\UW=6.9\times10^{5}\UW$。

第二部分：海水的焦耳热功率

对每个直线推进器，根据电阻定律：

$R=\rho\frac{L}{S}$，

代入数据得：$R=\rho\frac{L}{S}=\rho\frac{c}{ab}=0.2\times\frac{0.3}{0.3\times0.4}\UO=0.5\UO$。

由热功率公式，$P=I^{2}R$，

代入数据得：$P_{\text{单}}=I^{2}R=5.0\times10^{5}\UW$，

$P_{2}=12\times5.0\times10^{5}\UW=6\times10^{6}\UW$。

第三部分：单位时间内海水动能的增加量

设 $\Delta t$ 时间内喷出海水的质量为 $m$，

$P_{3}=12\times\frac{\Delta E_{k}}{\Delta t}$，

考虑到海水的初动能为零，

$\Delta E_{k}=E_{k}=\frac{1}{2}mv_{\text{水对地}}^{2}$，

$m=\rho_{m}bcv_{\text{水对地}}\Delta t$，

$P_{3}=12\times\frac{\Delta E_{k}}{\Delta t}=12\times\frac{1}{2}\rho_{m}bcv_{\text{水对地}}^{2}=4.6\times10^{4}\UW$。
}

\end{enumerate}

\end{document}
